% =========================
% report/reproducibility.tex
% =========================

\subsection{Environment Specification}
To guarantee that the experiments are fully reproducible, the following environment details should be recorded and included with the results:
\begin{itemize}
  \item \textbf{Operating System (OS):} The version of the OS used during the experiment (e.g., Ubuntu 20.04, macOS 10.15).
  \item \textbf{Hardware:} The CPU model, RAM size, and GPU (if used) details.
  \item \textbf{Python version:} The version of Python used for the implementation.
  \item \textbf{OpenCV version and build:} The specific version of OpenCV (including CPU/GPU configuration) used for feature detection and matching.
  \item \textbf{Dependency lock file:} A lock file (e.g., \texttt{requirements.txt} or \texttt{environment.yml}) to capture all necessary Python dependencies.
\end{itemize}

\subsection{Determinism Controls}
Even well-established computer vision pipelines can show small nondeterminism due to threading and internal heuristics. To control for this, the following actions should be taken:
\begin{itemize}
  \item Set fixed random seeds in Python, NumPy, and PyTorch to ensure that the random elements of the pipeline do not introduce variance.
  \item Set OpenCV’s thread count to a fixed value (e.g., \texttt{cv2.setNumThreads(1)}) to ensure deterministic behavior.
  \item Log all parameter values and configurations to a \texttt{config.json} file, which can be referenced for reproducibility.
\end{itemize}

\subsection{Reproducible Command Lines}
To allow others to reproduce the results exactly, the following command lines should be provided. These commands include paths to data, output directories, and configuration files:
% \begin{itemize}
%   % \item \texttt{python src/align.py --data_path ../data --config config.yaml --output ../results}
%   % \item \texttt{python src/evaluate.py --input ../results/overlays --metrics ../results/metrics}
% \end{itemize}

The complete reproducibility includes the exact dataset paths, configuration files, and thresholds used during evaluation.

\subsection{File-Based Output Contract}
For reproducibility and ease of use, each experimental run should generate specific output files:
\begin{itemize}
  \item \texttt{config.json:} Contains all configuration parameters used in the run, including preprocessing settings, model hyperparameters, and evaluation thresholds.
  \item \texttt{curves/}: Contains CSV files of the training and validation loss curves and other relevant metrics (e.g., MRE, success rate).
  \item \texttt{samples/}: Contains generated samples (e.g., overlay images, augmented data).
  \item \texttt{checkpoints/}: Contains model weights and optimizer state files.
  \item \texttt{metadata.json:} Contains device details (e.g., GPU model), library versions, and dataset split hashes to ensure full reproducibility.
\end{itemize}
These files provide a complete record of the experimental process, ensuring that the results are fully reproducible.

